% Source: ACM_Templates.pdf, PDF pp. 83-85; user virtual-tree preprocessing.
% Supplement: https://cp-algorithms.com/graph/lca_binary_lifting.html
\subsection{树上倍增、LCA 与路径距离}
\label{tree:lca}
\AlgorithmInfo
    {查询最近公共祖先、第 $k$ 级祖先、祖先关系、路径距离；同时提供虚树所需 DFS 序。}
    {预处理 $O(n\log n)$，LCA 与跳祖先 $O(\log n)$，祖先判定 $O(1)$；空间 $O(n\log n)$。}
    {通常输入一棵树；森林须提供每个连通块的根。根深度为 1，0 为不存在；所有路径权和绝对值小于 $4\times10^{18}$。}
\subsubsection*{原理}
记 $fa_{u,j}$ 为 $u$ 的 $2^j$ 级祖先，则
\[
    fa_{u,j}=fa_{fa_{u,j-1},j-1}.
\]
先把较深点跳到同深度；若尚不相同，就从大到小同时跳，要求两点的目标祖先不同。最后两点的父亲相同，便是 LCA。表高按点数动态确定，不固定写死 20 层。

以根到点的路径权和 $dis$ 为前缀，令 $p=LCA(u,v)$，则
\[
    dist(u,v)=dis_u+dis_v-2dis_p.
\]
树上路径唯一，该公式允许负边权，但不再代表一般图的最短路。\InlineCode{minn[u]} 另外记录根到 $u$ 路径的最小边权，根处为 \InlineCode{INF}，用于虚树最小割示例。
\subsubsection*{使用教程}
只需 LCA 或还需拆分路径时，也可使用第~\ref{tree:hld-lca}~节的树剖求法；它无需倍增表，预处理与空间均为 $O(n)$。

\InlineCode{init(n)} 后 \InlineCode{add(u,v,w)} 加无向边，每边一次，省略权值取 1；\InlineCode{build(root)} 选根预处理，省略根取 1。\InlineCode{jump(u,k)} 的 $k=0$ 返回自身，负数或越过根返回 0；\InlineCode{lca(u,v)} 返回 LCA，\InlineCode{ancestor(u,v)} 判断 $u$ 是否为 $v$ 的祖先（含自身），\InlineCode{distance(u,v)} 返回两点路径权和。换根重新构建，新数据重新初始化并加边。

只求 LCA 时用 \InlineCode{add(u,v)} 即可；$w$ 只参与路径距离和根路径最小边。默认权为 1 时，\InlineCode{distance} 返回路径边数。
\begin{lstlisting}[style=mtycpp,numbers=none]
TreeLCA g;
g.init(5);
g.add(1,2,3),g.add(1,3,2),g.add(2,4,4),g.add(2,5,1);
g.build(1);
int p=g.lca(4,5),ancestor=g.jump(4,2);
long long dist=g.distance(3,4);
bool inside=g.ancestor(2,5);
\end{lstlisting}
\InlineCode{dfn,out,rk,sz} 的含义见第~\ref{math:tree-basics}~节。对于无环森林，调用 \InlineCode{build(vector<int>{...})}，为每棵树各选一个根（如两棵树的根为 1、4，则传 \InlineCode{vector<int>{1,4}}）；单棵树直接传一个整数。根列表必须覆盖所有块，不能把含环图作为森林输入。跨块 LCA 为 0，距离为 \InlineCode{INF}。后续差分、长链与虚树封装均按连通树使用。
\needspace{7\baselineskip}
\subsubsection*{模板代码}
\CodeFile{code/trees/lca.cpp}
